\section{Skrótowa analiza wymagań}

Cel projektu: opracowanie metody (lub zestawu metod) pozwalającej odróżnić tekst napisany przez człowieka od tekstu wygenerowanego przez AI, z uwzględnieniem ograniczeń wynikających z konsultacji eksperckich oraz z obserwacji literatury (krótkie teksty, zmiana domeny, brak współpracy generatora).

\subsection{Wymagania funkcjonalne}

\begin{table}[H]
\centering
\small
\begin{tabularx}{\textwidth}{|c|X|c|}
\hline
\textbf{ID} & \textbf{Wymaganie} & \textbf{Priorytet} \\
\hline
W1 & Klasyfikacja binarna: tekst ludzki / tekst AI & wysoki \\
\hline
W2 & Minimalna dokładność klasyfikacji $\geq 51\%$ na wybranym zbiorze testowym & wysoki \\
\hline
W3 & Obsługa tekstów o ustalonej minimalnej długości (do ustalenia eksperymentalnie) & średni \\
\hline
W4 & Porównanie co najmniej jednego podejścia statystycznego i jednego ML (o ile pozwolą dane) & średni \\
\hline
W5 & Wykorzystanie otwartego zbioru danych (lub własnego podzbioru) & wysoki \\
\hline
W6 & Udokumentowanie ograniczeń metody (krótkie teksty, domena, język) & średni \\
\hline
W7 & Metoda działa bez dostępu do logitów generatora i bez znaku wodnego po stronie modelu (black-box / post-hoc) & wysoki \\
\hline
W8 & Raportowanie podstawowych metryk poza dokładnością (precyzja, czułość lub F1) & średni \\
\hline
\end{tabularx}
\caption{Lista wymagań funkcjonalnych}
\end{table}

\subsection{Wymagania niefunkcjonalne}

\begin{table}[H]
\centering
\small
\begin{tabularx}{\textwidth}{|c|X|c|}
\hline
\textbf{ID} & \textbf{Wymaganie} & \textbf{Priorytet} \\
\hline
N1 & Reprodukowalność eksperymentów (opis podziału danych, seed, wersje bibliotek) & wysoki \\
\hline
N2 & Możliwość uruchomienia baseline'u na sprzęcie studenckim (CPU lub jeden GPU) & średni \\
\hline
N3 & Kod i dokumentacja w repozytorium projektu & średni \\
\hline
N4 & Świadome rozróżnienie detekcji post-hoc od detekcji znaku wodnego dostawcy (np.\ textGrain) & średni \\
\hline
\end{tabularx}
\caption{Lista wymagań niefunkcjonalnych}
\end{table}

Poprzeczka dokładności została świadomie ustawiona nisko ($51\%+$), ponieważ na etapie planowania nie jest pewne, czy stabilna detekcja jest możliwa przy dostępnych zasobach i danych - zwłaszcza dla języka polskiego. Wymaganie W7 wynika z faktu, że komercyjne znakowanie wodne (np.\ OpenAI textGrain) obejmuje wyłącznie wybrane modele jednego dostawcy i nie zastępuje niezależnej detekcji tekstów z innych źródeł.
